\section{Bridge III: Nonlinear Compactness}

The second bridge upgrades compactness to the level required for nonlinear
closure. Once the cascade channel is controlled, the route still has to justify
passage to the limit in the convective term without importing regularized
structure. On the theorem surface this is achieved on one fixed classical
approximation scheme. The compactness input is the standard Aubin--Lions--Simon
framework on bounded sets together with the admissible energy-tail bound that
globalizes the \(\mathbb R^3\) limit on the declared class
\cite{Simon1987,BoyerFabrie2013,Temam2001}.

\begin{proposition}[Classical nonlinear passage]
\label{prop:compactness}
There exists a classical approximation scheme \(u^{(n)}\to u\) strongly in \(L^2([0,T];L^2_{\mathrm{loc}}(\mathbb{R}^3))\) such that
\[
u^{(n)}\otimes u^{(n)} \to u\otimes u
\quad\text{in distributions on } (0,T)\times\mathbb{R}^3.
\]
Hence the convective term passes to the limit on the classical equation with no additional nonlinear defect term.
\end{proposition}

\begin{proof}
The approximation family carries the standard uniform bounds
\[
u^{(n)} \text{ bounded in } L^\infty(0,T;L^2_\sigma(\mathbb R^3))
\cap L^2(0,T;H^1_\sigma(\mathbb R^3)),
\]
with \(\partial_t u^{(n)}\) bounded in \(L^2(0,T;H^{-1}_\sigma)\). Hence
Simon/Aubin--Lions gives, after extraction, strong convergence in
\(L^2(0,T;L^2(B_R))\) on each fixed ball \(B_R\)
\cite{Simon1987,BoyerFabrie2013}. The route's admissible energy-tail bound then
supplies the globalization step: for every \(\varepsilon>0\) there is
\(R<\infty\) such that
\[
\sup_n \int_0^T\!\!\int_{|x|>R} |u^{(n)}(t,x)|^2\,dx\,dt \le \varepsilon.
\]
By weak lower semicontinuity, the same tail bound holds for the limit \(u\),
and therefore
\[
\|u^{(n)}-u\|_{L^2_tL^2_x(\mathbb R^3)}
\le
\|u^{(n)}-u\|_{L^2_tL^2(B_R)} + 2\varepsilon^{1/2}.
\]
This yields strong convergence in \(L^2([0,T];L^2(\mathbb R^3))\) on the
declared theorem surface. Let
\(\varphi\in C_c^\infty((0,T)\times\mathbb{R}^3;\mathbb{R}^3)\) be
divergence-free. By the fixed classical approximation scheme from the
proposition,
\[
u^{(n)}\otimes u^{(n)} \to u\otimes u
\quad\text{in distributions on } (0,T)\times\mathbb{R}^3.
\]
Therefore
\[
\int_0^T\!\!\int_{\mathbb{R}^3}
\bigl(u^{(n)}\otimes u^{(n)}\bigr):\nabla\varphi\,dx\,dt
\to
\int_0^T\!\!\int_{\mathbb{R}^3}
\bigl(u\otimes u\bigr):\nabla\varphi\,dx\,dt.
\]
The strong local \(L^2\) convergence identifies the limiting field, while the tensor convergence excludes any extra defect contribution at the theorem level. This is exactly the nonlinear passage needed in the present route.
\end{proof}

The compactness step is therefore no longer a placeholder. It is the fixed
nonlinear-closure package imported by the main theorem, and on
\(\mathbb R^3\) its local-to-global passage is tied explicitly to the same
admissible energy-tail hypothesis already carried by the theorem surface.
